\begin{frame}{각 요소 하나씩 (2/2): $P$, $R$, $\gamma$}
  \kb{$P$, 전이확률}
  \begin{itemize}
    \item $P(s'|s,a)$ 는 ``이 게임의 \textbf{물리 법칙}''이다.
    \item Gymnasium 환경에서는 이 함수가 \texttt{step()} 함수의
      내부 구현으로 감춰져 있고, 우리는 \textbf{샘플}(실제 전이)만
      관찰할 수 있다.
    \item \textbf{모델 기반}(Week 4--5 동적계획법): $P$ 를 (추정해서)
      명시적으로 써서 계획. \textbf{모델 없는}(Week 6--11): $P$ 를
      전혀 쓰지 않고 경험(실제 전이 샘플)에서만 가치를 배운다.
    \item 같은 $P$ 라도 ``정식적 표''로 줄 것이냐 ``샘플 공급
      원''으로 둘 것이냐가 알고리즘의 가족을 갈라놓는
      \textbf{첫 번째 갈림길}.
  \end{itemize}
  \vskip0.4em
  \kb{$R$, 보상}
  \begin{itemize}
    \item 수학적으로 더 정확한 표기는 \textbf{전이} $R(s,a,s')$ 에
      붙는 것(같은 $(s,a)$에서도 다음 상태가 다르면 보상이 다를
      수 있음 --- ``도착했다/아직이다''는 $s'$ 를 봐야 알 수).
    \item 보상의 \textbf{단위$\cdot$스케일}은 알고리즘에 직접
      영향을 준다: 10배 키우면 가치함수 전체가 10배가 되지만
      \textbf{최적 정책은 변하지 않는다}(연습문제 5).
    \item 그러나 \textbf{상태/행동 사이에서 불균일}하게 들쭉한 스케일
      은 조심해야 한다(Week 5에서 탐험 균형이 깨지는 사례).
  \end{itemize}
  \vskip0.3em
  {\footnotesize \kb{$\gamma$, 할인율}: $\gamma < 1$ 은 ``에이전트가
  영원히 살지 않는 세계''를 뜻하며, $V^\pi, Q^\pi$ 가 \textbf{유계
  함수}로 존재하게 하는 수학적 필요 조건. (역할은 3.2에서 상세히)}
\end{frame}
